\itemhead{p.82 练习 · 第2题（教材p.82）}{XJ-B1-C03-PR-P082-Y1193}
\label{XJ-B1-C03-PR-P082-Y1193}
{\footnotesize 湘教版教材p.82 / PDF第89页 / p.82 练习 / 第2题}
\field{原题}{检验下列函数的增减性，并说明是否有最大（小）值。如果有，指出最大（小）值和对应的最大（小）值点。
\subq{1}$f(x)=-\dfrac2x\ (x\in(-\infty,0))$；
\subq{2}$f(x)=3-\dfrac x2\ (x\in[-6,1])$；
\subq{3}$f(x)=x^2-6x+7\ (x\in[-2,4])$；
\subq{4}$f(x)=\dfrac x{1+x}\ (x\in[0,3])$。}
\field{参考答案}{（1）递增，无最大最小值；（2）递减，最大值6在 $x=-6$，最小值 $5/2$ 在 $x=1$；（3）在 $[-2,3]$ 递减、$[3,4]$ 递增，最小值 $-2$ 在3，最大值23在 $-2$；（4）递增，最小值0在0，最大值 $3/4$ 在3。}
{\small \field{关联知识}{F05 单调性定义、证明与参数（教材帮p.97、98、100 / PDF106、107、109）；F06 最大值、最小值与值域（教材帮p.99、83、84 / PDF108、92、93）}
\field{方法与易错点}{用定义比较任意两值，或依据已知图象性质；含参时使指定区间被相应单调区间覆盖。 易错：用有限采样代替证明；把分离半轴合并；漏跨零点情形。}
\field{答案依据}{答案册相关题组：PDF第22、23页（册内页21、22）；按当前题库原题逐项重算，不按旧版题号机械对套。}}
